
\documentclass[11pt]{article}

\usepackage[margin=1in]{geometry}
\usepackage{amsmath,amssymb,amsthm,mathtools}
\usepackage{enumitem}
\usepackage[hidelinks]{hyperref}
\usepackage[T1]{fontenc}
\usepackage{lmodern}

\newtheorem{theorem}{Theorem}
\newtheorem{lemma}{Lemma}
\newtheorem{remark}{Remark}

\newcommand{\OO}{\mathcal O}
\newcommand{\Bl}{\operatorname{Bl}}
\newcommand{\Spec}{\operatorname{Spec}}
\newcommand{\Proj}{\operatorname{Proj}}

\title{Finite-Degree Approximation on the Blow-Up Itself}
\author{}
\date{}

\begin{document}

\maketitle

\begin{theorem}[Finite-degree approximation on the blow-up itself]
Let $X$ be a quasi-compact and quasi-separated scheme, let $I$ be a finite-type
quasi-coherent ideal sheaf on $X$, and let
\[
  \pi:Y=\Bl_I(X)\longrightarrow X
\]
be the blow-up of $X$ along $I$.
Let $J$ be a finite-type quasi-coherent ideal sheaf on $Y$, and write
\[
  I_E := I\OO_Y = \pi^{-1}I\cdot \OO_Y
\]
for the exceptional ideal.

Then there exist an integer $d\geq 1$ and a finite-type quasi-coherent ideal
sheaf $L\subset I^d$ such that
\[
  L\OO_Y = I_E^d J.
\]
Equivalently, in the notation of the Lean statement,
\[
  L.\operatorname{comap}(I.\operatorname{blowupMap})
  =
  \bigl(I.\operatorname{comap}(I.\operatorname{blowupMap})\bigr)^d * J.
\]
\end{theorem}

\begin{proof}
Let
\[
  \pi:Y=\Bl_I(X)\longrightarrow X
\]
be the blow-up, and write
\[
  I_E := I\OO_Y
\]
for the exceptional ideal.
We want to construct $d>0$ and a finite-type ideal $L\subset I^d$ whose
inverse image to $Y$ is exactly $I_E^dJ$.

The proof is organized around the following maximal degree-$d$ candidate:
\[
  C_d := I^d \cap \pi_*(I_E^dJ).
\]
More precisely, $C_d$ is the subsheaf of $I^d$ consisting of sections whose
pullback to the blow-up lies in $I_E^dJ$.
Thus $C_d$ is the largest possible degree-$d$ ideal whose inverse image is
contained in $I_E^dJ$.

The main point is to show that for all sufficiently large $d$,
\[
  C_d\OO_Y = I_E^dJ.
\]
After that, since $C_d$ need not itself be of finite type, we will select
finitely many local generators inside $C_d$ and place them in a finite-type
quasi-coherent subideal $L\subset C_d$.

\medskip

\noindent
\textbf{Step 1: Choose finitely many affine blow-up charts.}

Since $X$ is quasi-compact, choose a finite affine open cover
\[
  X=\bigcup_{\alpha=1}^N V_\alpha,
  \qquad
  V_\alpha=\Spec(A_\alpha).
\]
Because $I$ is of finite type, after refining this cover if necessary we may
write
\[
  I(V_\alpha)
  =
  I_\alpha
  =
  (a_{\alpha,1},\dots,a_{\alpha,r_\alpha}).
\]

Over $V_\alpha$, the blow-up is
\[
  \pi^{-1}(V_\alpha)
  =
  \Proj\!\left(\bigoplus_{n\geq 0} I_\alpha^n\right),
\]
and it is covered by the standard affine blow-up charts
\[
  U_{\alpha,i}
  =
  \Spec(B_{\alpha,i}),
  \qquad
  B_{\alpha,i}
  =
  A_\alpha\!\left[\frac{I_\alpha}{a_{\alpha,i}}\right].
\]

Fix $\alpha$ temporarily and suppress it from the notation. Thus we write
\[
  A,\qquad
  I=(a_1,\dots,a_r),\qquad
  B_i=A[I/a_i],\qquad
  U_i=\Spec(B_i).
\]

On the chart $U_i$, the pullback of $I$ is principal:
\[
  I_EB_i = IB_i = a_iB_i.
\]
Indeed, for every $b\in I$,
\[
  b = a_i\cdot \frac{b}{a_i}
\]
in $B_i$.
Consequently,
\[
  (I_E^dJ)|_{U_i}
  =
  a_i^dJ_i,
\]
where
\[
  J_i := J(U_i)\subset B_i.
\]

Since $J$ is finite type and $U_i$ is affine, choose finitely many generators
\[
  \mu_{i,1},\dots,\mu_{i,s_i}
\]
for $J_i$.

Thus, in order to prove that $C_d\OO_Y$ contains $I_E^dJ$, it is enough to
produce, for each generator $\mu_{i,k}$ and for one common sufficiently large
$d$, an element
\[
  m_{i,k}\in C_d(V)
\]
such that
\[
  \frac{m_{i,k}}{a_i^d}
  =
  \mu_{i,k}
\]
on $U_i$.
Indeed, this equality is equivalent to
\[
  m_{i,k}=a_i^d\mu_{i,k}
\]
as a function on $U_i$, and these elements generate
\[
  a_i^dJ_i=(I_E^dJ)|_{U_i}.
\]

\medskip

\noindent
\textbf{Step 2: The same-Rees-algebra denominator-clearing argument.}

Fix one chart $U_i$ and one generator
\[
  \mu\in J_i.
\]
Because
\[
  B_i=A[I/a_i]
\]
is the degree-zero part of the localization of the Rees algebra, every element
of $B_i$ may be represented in the form
\[
  \mu=\frac{g}{a_i^e}
\]
for some $e\geq 0$ and some
\[
  g\in I^e.
\]

If there were no other charts, then for any $d\geq e$ we could simply put
\[
  m=a_i^{d-e}g.
\]
Then $m\in I^d$ and
\[
  \frac{m}{a_i^d}
  =
  \frac{a_i^{d-e}g}{a_i^d}
  =
  \frac{g}{a_i^e}
  =
  \mu.
\]
The only issue is to ensure that the same numerator $m$ satisfies the
required membership condition on every other blow-up chart.

Fix another chart $U_j$.
Inside
\[
  B_j=A[I/a_j]
\]
consider the genuine element
\[
  x_j:=\frac{g}{a_j^e}.
\]
This is an element of $B_j$ because $g\in I^e$.

On the overlap $U_i\cap U_j$, we have
\[
  x_j
  =
  \frac{g}{a_j^e}
  =
  \left(\frac{a_i}{a_j}\right)^e
  \frac{g}{a_i^e}
  =
  \left(\frac{a_i}{a_j}\right)^e\mu.
\]
Since $\mu$ is a section of $J$ on $U_i$, its restriction to the overlap lies
in the restriction of $J$ there.
Therefore $x_j$ lies in the localization of $J_j$ obtained by inverting
\[
  s_{ij}:=\frac{a_i}{a_j}\in B_j.
\]
Equivalently,
\[
  x_j\in (J_j)_{s_{ij}}.
\]

By the elementary localization criterion for submodule membership, there
exists an integer $N_j\geq 0$ such that
\[
  s_{ij}^{N_j}x_j\in J_j.
\]
In other words,
\[
  \left(\frac{a_i}{a_j}\right)^{N_j}
  \frac{g}{a_j^e}
  \in J_j.
\]

There are only finitely many charts $U_j$, so choose
\[
  N\geq N_j
\]
for all $j$.
Because $J_j$ is an ideal, this implies
\[
  \left(\frac{a_i}{a_j}\right)^N
  \frac{g}{a_j^e}
  \in J_j
\]
for every $j$.

Now put
\[
  d_0:=e+N,
  \qquad
  m:=a_i^Ng.
\]
Then
\[
  m\in I^{e+N}=I^{d_0}.
\]

On the $i$-th chart,
\[
  \frac{m}{a_i^{d_0}}
  =
  \frac{a_i^Ng}{a_i^{e+N}}
  =
  \frac{g}{a_i^e}
  =
  \mu.
\]

On the $j$-th chart,
\[
  \frac{m}{a_j^{d_0}}
  =
  \frac{a_i^Ng}{a_j^{e+N}}
  =
  \left(\frac{a_i}{a_j}\right)^N
  \frac{g}{a_j^e}
  \in J_j.
\]
Therefore
\[
  m
  =
  a_j^{d_0}
  \left(\frac{m}{a_j^{d_0}}\right)
  \in
  a_j^{d_0}J_j
  =
  (I_E^{d_0}J)|_{U_j}.
\]
Since the $U_j$ cover $\pi^{-1}(V)$, the pullback of $m$ belongs to
$I_E^{d_0}J$ everywhere over $V$.
Hence
\[
  m\in C_{d_0}(V).
\]
Moreover,
\[
  \frac{m}{a_i^{d_0}}=\mu.
\]

This proves the desired construction for one chosen local generator.

\medskip

\noindent
\textbf{Step 3: The construction persists in all larger degrees.}

Suppose
\[
  m\in C_{d_0}(V)
\]
and
\[
  \frac{m}{a_i^{d_0}}=\mu.
\]
For any $d\geq d_0$, define
\[
  m^{(d)}
  :=
  a_i^{d-d_0}m.
\]
Then $m^{(d)}\in I^d$, and on $U_i$,
\[
  \frac{m^{(d)}}{a_i^d}
  =
  \frac{m}{a_i^{d_0}}
  =
  \mu.
\]

On another chart $U_j$,
\[
  \frac{m^{(d)}}{a_j^d}
  =
  \left(\frac{a_i}{a_j}\right)^{d-d_0}
  \frac{m}{a_j^{d_0}}.
\]
Since $m\in C_{d_0}(V)$, the second factor lies in $J_j$, and the first factor
lies in $B_j$.
Because $J_j$ is an ideal,
\[
  \frac{m^{(d)}}{a_j^d}\in J_j.
\]
Thus
\[
  m^{(d)}\in C_d(V).
\]

Therefore each chosen local generator of $J$ can be represented by an element
of $C_d$ for every sufficiently large $d$.

\medskip

\noindent
\textbf{Step 4: Choose one global degree.}

Return to the finite affine cover
\[
  X=\bigcup_\alpha V_\alpha.
\]
For every triple
\[
  (\alpha,i,k)
\]
consisting of an affine open $V_\alpha$, a blow-up chart
$U_{\alpha,i}$, and a chosen generator $\mu_{\alpha,i,k}$ of $J$ on that
chart, the previous argument produces a threshold
\[
  d_{\alpha,i,k}
\]
such that for every $d\geq d_{\alpha,i,k}$ there exists
\[
  m_{\alpha,i,k}\in C_d(V_\alpha)
\]
with
\[
  \frac{m_{\alpha,i,k}}{a_{\alpha,i}^d}
  =
  \mu_{\alpha,i,k}.
\]

There are only finitely many such triples.
Choose one integer
\[
  d\geq 1
\]
larger than all the $d_{\alpha,i,k}$.

Then for every $(\alpha,i,k)$ we have a section
\[
  m_{\alpha,i,k}\in C_d(V_\alpha)
\]
such that
\[
  m_{\alpha,i,k}
  =
  a_{\alpha,i}^d\mu_{\alpha,i,k}
\]
on $U_{\alpha,i}$.

Since the $\mu_{\alpha,i,k}$ generate $J$ on $U_{\alpha,i}$, the sections
$m_{\alpha,i,k}$ generate
\[
  a_{\alpha,i}^dJ
  =
  I_E^dJ
\]
there.
Thus the inverse image of $C_d$ contains $I_E^dJ$ locally on every blow-up
chart.

The reverse inclusion is built into the definition of $C_d$: every section
of $C_d$ is a section of $I^d$ whose pullback lies in $I_E^dJ$.
Hence
\[
  C_d\OO_Y=I_E^dJ.
\]

\medskip

\noindent
\textbf{Step 5: Extract a finite-type subideal.}

The ideal $C_d$ need not itself be of finite type.
However, only finitely many local sections
\[
  m_{\alpha,i,k}
\]
are needed to generate $I_E^dJ$ after pullback.

Because $X$ is quasi-compact and quasi-separated, a quasi-coherent module is
the directed union of its finite-type quasi-coherent submodules.
Therefore there exists a finite-type quasi-coherent submodule
\[
  L\subset C_d
\]
containing all the finitely many sections $m_{\alpha,i,k}$.

Since
\[
  C_d\subset I^d\subset \OO_X,
\]
the sheaf $L$ is itself a quasi-coherent ideal sheaf, and
\[
  L\subset I^d.
\]

We now show that
\[
  L\OO_Y=I_E^dJ.
\]

Because $L\subset C_d$, the defining property of $C_d$ gives
\[
  L\OO_Y\subset I_E^dJ.
\]

For the reverse inclusion, work on one chart $U_{\alpha,i}$.
There,
\[
  I_E^dJ
  =
  a_{\alpha,i}^dJ.
\]
The ideal $J$ is generated by the finitely many
\[
  \mu_{\alpha,i,1},\dots,\mu_{\alpha,i,s_i},
\]
so $I_E^dJ$ is generated by
\[
  a_{\alpha,i}^d\mu_{\alpha,i,k}.
\]
By construction,
\[
  a_{\alpha,i}^d\mu_{\alpha,i,k}
  =
  m_{\alpha,i,k},
\]
and each $m_{\alpha,i,k}$ lies in $L(V_\alpha)$.
Therefore every generator of $I_E^dJ$ on $U_{\alpha,i}$ belongs to
$L\OO_Y$.
Hence
\[
  I_E^dJ|_{U_{\alpha,i}}
  \subset
  L\OO_Y|_{U_{\alpha,i}}.
\]

Since the charts $U_{\alpha,i}$ cover $Y$, this gives globally
\[
  I_E^dJ\subset L\OO_Y.
\]
Combining the two inclusions,
\[
  L\OO_Y=I_E^dJ.
\]

Thus we have constructed an integer $d>0$ and a finite-type quasi-coherent
ideal sheaf $L$ satisfying
\[
  L\subset I^d
  \qquad\text{and}\qquad
  L\OO_Y=I_E^dJ.
\]
This is exactly the desired conclusion.
\end{proof}

\bigskip

\begin{remark}[Formalization structure]
For a Lean formalization, the argument naturally separates into the following
four auxiliary statements.

\begin{enumerate}[label=\arabic*.]
  \item
  \emph{Blow-up chart description:}
  over $V=\Spec(A)$ with $I=(a_1,\dots,a_r)$, the blow-up is covered by
  $\Spec(A[I/a_i])$, and on the $a_i$-chart the exceptional ideal is
  generated by $a_i$.

  \item
  \emph{Same-Rees-algebra clearing lemma:}
  if
  \[
    \mu=\frac{g}{a_i^e}\in J_i,
    \qquad g\in I^e,
  \]
  then for every $j$ one has
  \[
    \frac{g}{a_j^e}\in B_j,
  \]
  and after multiplying by a sufficiently large power of $a_i/a_j$ the
  result lies in $J_j$.
  Choosing one exponent valid for every $j$ gives
  \[
    m=a_i^Ng,\qquad d=e+N,
  \]
  with
  \[
    \frac{m}{a_i^d}=\mu,
    \qquad
    \frac{m}{a_j^d}\in J_j
    \quad\text{for all }j.
  \]

  \item
  \emph{Uniform-degree lemma:}
  once a chosen generator is represented in degree $d_0$, multiplying the
  numerator by $a_i^{d-d_0}$ produces a representation in every degree
  $d\geq d_0$.
  Finiteness of the affine cover, blow-up charts, and generator sets then
  gives a single global degree $d$.

  \item
  \emph{Finite-type extraction:}
  place all selected local numerators inside the maximal candidate
  \[
    C_d
    =
    I^d\cap \pi_*(I_E^dJ).
  \]
  The qcqs hypothesis gives a finite-type quasi-coherent subideal
  \[
    L\subset C_d
  \]
  containing them.
  The inclusion $L\subset C_d$ gives one inclusion after pullback, and the
  selected chart generators give the reverse inclusion.
\end{enumerate}

Thus the only genuinely nonformal algebraic step is the denominator-clearing
argument on two affine charts belonging to the same Rees algebra.
\end{remark}

\end{document}
